\selectlanguage{ngerman}%
\begin{ledgroupsized}[r]{120mm}%
\footnotesize%
\pstart%
\noindent%
\textbf{Überlieferung:}%
\pend%
\end{ledgroupsized}%
\begin{ledgroupsized}[r]{114mm}%
\footnotesize%
\pstart%
\parindent -6mm%
\makebox[6mm][l]{\textit{L}}%
Aufzeichnung:
LH~XXXVII~5 Bl.\;196\textendash197.
Ein Bogen~2\textsuperscript{o},
Ränder ausgefranst mit geringfügigem Textverlust am unteren Rand von Bl.\;196\,r\textsuperscript{o};
Papiererhaltungsmaßnahmen.
Ein \href{\Scan{LH_37_05_196wz}}{Wasserzeichen} auf Bl.\;196 %
mit \href{\Scan{LH_37_05_197wz}}{Gegenmarke} auf Bl.\;197
Drei Seiten auf Bl.\;196\,r\textsuperscript{o}\! bis Bl.\;197\,r\textsuperscript{o}\! und zwei Zeilen auf Bl.\;197\,v\textsuperscript{o}\!
(sonst %
unbeschrieben). %
\pend
\end{ledgroupsized}
\vspace{5mm}%
\begin{ledgroup}%
\footnotesize%
\pstart%
\noindent%
\textbf{Datierungsgründe:}
Die Datierung der Aufzeichnung N.\;\ref{60359} beruht auf äußeren Eigenschaften des Text\-trägers %
sowie auf inhaltlichen Betrachtungen.
Siehe die ausführliche Begründung auf S.\;\pageref{60359_dat}.%
\pend%
\end{ledgroup}%
\selectlanguage{latin}%
\frenchspacing%
\vspace{8mm}
\pstart%
\normalsize%
\noindent%
\lbrack%
\href{\Scan{LH_37_05_196r}}%
{196\,r\textsuperscript{o}}%
\rbrack\ %
Videamus
qua ratione%
\indexS{ratio diminutionis}
diminuatur corpus incurrens%
\indexS{corpus incurrens in aliud quiescens}
semper in alia quiescentia%
\indexS{corpus quiescens ante ictum}
\edtext{nova,}{%
\lemma{}%
\Bfootnote{%
nova%
~\textit{erg.\:L}%
}}
ita ut cuilibet eorum
eandem semper tribuat celeritatem.%
\indexS{celeritas tributa}
Si corpus \textit{a}
\edtext{%
celeritate~\textit{v}%
}{%
\lemma{}%
\Bfootnote{%
celeritate \textit{v}%
~\textit{erg.\:L}%
}}
incurrat in corpus \textit{b} %
celeritate \textit{y}, %
et post ictum%
\indexS{ictus}
illud habeat celeritatem \textit{x},
hoc celeritatem~\textit{z},%
\indexS{celeritas post ictum}
ponaturque ambo ire ad easdem partes
ante ictum et post ictum,%
\indexS{corpus iens ad easdem partes ante et post ictum}
fiet
$v + x = y + z$
et
$av + by = ax + bz,$
fietque
$x = y + z - v = av + by - bz,\, :\, a,$
ergo fiet
$ay + az - av = av + by - bz,$
adeoque fiet
$z = 2av + (b - a)y ,\, :\: , a + b,$
et similiter
$x =
\edtext{\lbrack2by\rbrack}{\lemma{}\Bfootnote{%
2\textit{ay} \textit{L~ändert Hrsg.}}}
+ (a - b) v ,\, :\: , a + b.$
Quodsi \textit{b} ante ictum fuerit quiescens,%
\indexS{corpus quiescens ante ictum}
seu $y = 0,$
fiet $x = (a - b)v : (a + b)$
et
$z = 2av : (a + b),$
adeoque \textit{x} ad \textit{z}
ut $a - b$ ad 2\textit{a}.
\pend 
\vspace{1.0em}%
\pstart
\noindent
\lbrack%
\textit{Nachfolgend kleingedruckter Text in L gestrichen:}\rbrack %
\pend
\vspace{0.25em}
\footnotesize
\pstart%
\noindent%
Sit jam corpus \textit{c},
cui a corpore \textit{a}
habente
\edtext{%
velocitatem \textit{x},%
\indexS{velocitas corporis incurrentis}
quam et vocabimus \textlangle\textit{v},\textrangle\
danda sit velocitas%
}{%
\lemma{velocitatem \textit{x},}%
\Bfootnote{%
\textit{(1)}~vel \textit{(v)}
\textit{(2)}~quam et vocabimus
\textit{(a)}~\textit{(v)}
\textit{(b)}~\textlangle\textit{v},\textrangle\
\textit{(aa)}~velo
\textit{(bb)}~danda sit velocitas%
~\textit{L}%
}}
aequalis ipsi \textit{z},
utique velocitas%
\indexS{velocitas accepta}
\edtext{%
corporis \textit{c} debebit esse
$2ax : (a + c) = 2av : (a + b),$%
}{%
\lemma{corporis \textit{c}}%
\Bfootnote{%
\textit{(1)}~debebit esse 2\textit{av}
\textit{(2)}~debebit esse
\textit{(a)}~2\textit{av}
\textit{(b)}~$2ax : (a + c) = 2av : (a + b),$%
~\textit{L}%
}}
ergo fiet $a + c$
\edtext{%
ad $a + b,$
ut \textit{x} ad \textit{v}%
}{%
\lemma{ad $a + b,$ ut}%
\Bfootnote{%
\textit{(1)}~\textit{(v)}
\textit{(a)}~ad \textit{v}
\textit{(b)}~vel
\textit{(c)}~$\langle\sqcap\rangle$
\textit{(d)}~vel \textit{x}
\textit{(2)}~\textit{x} ad \textit{v}%
~\textit{L}%
}}
seu $a + c$ ad $a + b$
ut $(a - b)v : (a + b)$ ad \textit{v},
ergo $a + c$ ad $a + b$
ut $a - b$ ad $a + b.$
Unde sequeretur $c = -\,b,$
quod est absurdum.%
\indexS{absurdum}
Ergo impossibile est
reperiri corpus quiescens%
\indexS{corpus quiescens ante ictum}
cui corpus motum celeritate residua%
\indexS{celeritas residua post ictum}
tantum det celeritatis,
quantum idem corpus motum dederat antea
alteri quiescenti%
\indexS{corpus quiescens ante ictum}
celeritate integra%
\indexS{celeritas integra}%
\lbrack,\rbrack\
intelligendo incursus
\edlabel{37_05_196-197_1a}%
directos.%
\indexS{incursus directus}
\edtext{}{%
\xxref{37_05_196-197_1a}{37_05_196-197_1b}%
\lemma{directos.}%
\Bfootnote{%
\textit{(1)}~Vel quod idem est: imminuere corporis motu
\textit{(2)}~Hinc sequitur problema istud generaliter solvi non posse:
\textit{(3)}~Tentemus jam problema:%
~\textit{L}%
}}%
\pend
\vspace{1.0em}%
\normalsize %
\pstart%
\noindent%
Tentemus jam problema:%
\indexS{problema}%
\edlabel{37_05_196-197_1b}%
\textls{ %
Data corporis magnitudine%
\indexS{magnitudo corporis incurrentis}
et velocitate,%
\indexS{velocitas corporis incurrentis}
invenire magnitudinem corporis quiescentis talem,%
\indexS{magnitudo corporis quiescentis}
ut
\edtext{hoc}{%
\lemma{}%
\Bfootnote{%
\textls{hoc}%
~\textit{erg.\:L}%
}}
ab illo directe impulsum
velocitatem accipiat%
\indexS{velocitas accepta}
 }%
\edtext{%
\textls{datam.}
Sit datum corpus%
\indexS{corpus incurrens in aliud quiescens}%
}{%
\lemma{\textls{datam.}}%
\Bfootnote{%
\textit{(1)}~Tentemus hoc problema %
\textit{(2)}~Sit
\textit{(a)}~qui
\textit{(b)}~datum corpus%
~\textit{L}%
}}
motum \textit{a} %
cujus velocitas \textit{v},%
\indexS{velocitas corporis incurrentis}
sit quiescens~\textit{b},%
\indexS{corpus quiescens ante ictum}
quaeritur quantum debeat esse \textit{b},
ut velocitas hujus \textit{z} %
quam acciperet %
si ab ipso \textit{a} directe impelleretur%
\indexS{corpus impellens directe}%
\indexS{corpus impulsum directe}%
\lbrack,\rbrack\
esset data.%
\indexS{velocitas accepta}%
\indexS{velocitas data}
Itaque data erit
$z = 2av :\, , a + b.$
Ergo
$az + bz = 2av.$
Ergo fiet 
$b = %
\edtext{%
2av - az\, ,\, : z.$
Hinc%
}{%
\lemma{}%
\Bfootnote{%
$2av - az\, ,\, : z$
\bars{, seu $2av : z\, , -\, a$ \textit{gestr.}}%
. Hinc%
~\textit{L}%
}}
sequitur%
\lbrack:\rbrack\
ut problema solvi possit%
\indexS{problema}
necesse esse
ut $2v - z$ sit quantitas affirmativa,%
\indexS{quantitas affirmativa}
seu ut 2\textit{v} sit major quam \textit{z}.
Nempe \textit{b} debet esse ad \textit{a},
ut $2v - z$ ad \textit{z}.
\pend%
\pstart%
Ponamus jam corpus \textit{a}
\edtext{%
velocitate sua \textit{v}%
}{%
\lemma{}%
\Bfootnote{%
velocitate sua \textit{v}%
~\textit{erg.\:L}%
}}
corpori \textit{b}
\edtext{%
dedisse celeritatem \textit{z};
et videamus%
}{%
\lemma{dedisse}%
\Bfootnote{%
\textit{(1)}~motum \textit{z}. Videamus jam
\textit{(2)}~celeritatem \textit{z}; et videamus%
~\textit{L}%
}}
an liceat
\edtext{%
efficere
ut idem \textit{a}
residua velocitate \textit{(v)}
corpori alteri \textit{(b)}
det eandem velocitatem~\textit{z}.%
}{%
\lemma{efficere}%
\Bfootnote{%
\textit{(1)}~ut residua sua %
\textit{(2)}~ut idem \textit{a} residua velocitate
\textit{(a)}~\textit{(v)} corpori \textit{(b)} det velocitatem eand
\textit{(b)}~\textit{(v)} corpori %
\lbrack...\rbrack\ velocitatem~\textit{z}.%
~\textit{L}%
}}
Ergo \textit{(b)} debebit esse ad \textit{a},
ut $2\textit{(}v\textit{)} - z$ ad \textit{z}.
Ergo oportet ut 2\textit{(v)} sit major quam~\textit{z}.
Sed \textit{(v)},
id est \textit{x},
est
$a - b, v :\, ,a + b.$
Ergo
$2\,\overline{a - b}\, v :\, , a + b$
debet esse major quam \textit{z},
seu major quam
$2av :\, ,a + b,$
seu $a - b$ debet esse major
quam \textit{a}.
Quod est absurdum.%
\indexS{absurdum}
Itaque desideratum%
\indexS{desideratum}
praestari non
\edtext{%
potest,
si incursus sit directus%
\indexS{incursus directus}
et corpus%
}{%
\lemma{potest,}%
\Bfootnote{%
\textit{(1)}~si corpus %
\textit{(2)}~si incursus %
\lbrack...\rbrack\ et corpus%
~\textit{L}%
}}
excipiens%
\indexS{corpus excipiens}
sit minus,
adeoque incurrens post incursum pergat.%
\indexS{corpus incurrens in aliud quiescens}%
\indexS{corpus pergens post incursum}
\pend%
\pstart%
Sed quid
si \textit{b} sit
major
quam \textit{a},
tunc \textit{z} quidem tendet
ea via qua \textit{v},
sed \textit{a} reflectetur%
\indexS{corpus reflexum post ictum}
et \textit{x} erit quantitas negativa;%
\indexS{quantitas negativa}
sed cujus
\edtext{\lbrack magnitudinem\rbrack}{%
\lemma{}\Bfootnote{%
molem%
\indexS{moles}%
~\textit{L~ändert Hrsg.}}}
absolute considerando
fiet
$x = \overline{b - a}\,v : \overline{a + b}.$
Ipsum igitur \textit{a}
celeritate \textit{x} recurrens%
\indexS{corpus recurrens}%
\indexS{celeritas corporis recurrentis}
\edtext{%
vel \textit{(v)},%
}{%
\lemma{}%
\Bfootnote{%
vel \textit{(v)}%
~\textit{erg.\:L}%
}}
ponamus rursus incurrere%
\indexS{corpus incurrens in aliud quiescens}
in corpus \textit{(b)}
et dare ipsi \textit{(b)} celeritatem \textit{z}.
Ergo erit \textit{(b)} ad \textit{a}
ut $2\textit{(v)} - z$ ad \textit{z}.
Ergo debet esse 2\textit{(v)} seu 2\textit{x}
major quam \textit{z},
seu
$2\, \overline{b - a}\, v : \overline{a + b}$
major quam
$2av : \overline{a + b},$
seu erit
$b - a$
major
\edtext{%
quam \textit{a},
seu \textit{b} major quam 2\textit{a},
quod est%
}{%
\lemma{quam \textit{a},}%
\Bfootnote{%
\textit{(1)}~quod est %
\textit{(2)}~seu \textit{b} %
\lbrack...\rbrack\ quod est%
~\textit{L}%
}}
possibile.
Ergo \textit{(b)} ad \textit{a},
ut
\edtext{%
$2\, \overline{b - a}\, v : \overline{a + b}
-2av : \overline{a + b}$
ad $2av : \overline{a + b}\,;$%
}{%
\lemma{}%
\Bfootnote{%
$2\, \overline{b - a}\, v : \overline{a + b}$
\bars{$\!-\, 2av : \overline{a + b}$ \textit{erg.}}
ad $2av : \overline{a + b}\,;$%
~\textit{L}%
}}
seu \textit{(b)} ad \textit{a} %
ut $b - 2a$ ad~\textit{a}.
Ergo
$\textit{(b)} =
\edtext{%
b - 2a,$
et
\textit{(x)} vel
$\textit{((v))}%
}{%
\lemma{$b - 2a,$}%
\Bfootnote{%
\textit{(1)}~et \textit{(x)} %
\textit{(2)}~et \textit{(x)} vel \textit{((v))}%
~\textit{L}%
}}
= \overline{\textit{(b)} - a}\, \textit{(v)} : \overline{a + \!\textit{(b)}},$%
\rule[-1,0mm]{0mm}{5mm}
qua \textit{a} rursum reflectetur,%
\indexS{corpus reflexum post ictum}%
\indexS{celeritas corporis reflexi}
vel
$\textit{((v))} = \overline{b - 3a}\, %
\ovalbox{\overline{b - a}} v :
\edlabel{LH_37_05_196-197_196r-v_uebergang-1}%
\overline{a + b},\ovalbox{\overline{b - a}}.$%
\edtext{}{%
\xxref{LH_37_05_196-197_196r-v_uebergang-1}{LH_37_05_196-197_196r-v_uebergang-2}%
\lemma{$\overline{a + b}, \ovalbox{\overline{b - a}}.$}%
\Bfootnote{%
\textit{(1)}~Eodem modo %
$\textit{((b))} = \textit{(b)} - a = b - 2a,$ %
\textit{(a)}~et %
$\textit{(((b)))} = \textit{((b))} - a = b - \langle\--\rangle$ %
\textit{(b)}~et %
$\textit{(((b)))} = \textit{((b))} - a = b - 3a$ et %
$\textit{(((v)))} = \overline{\textit{(b)} - 2a},\, \overline{\textit{(b)} -a}\, (v) \,+\, \overline{a + \textit{(b)},\, \textit{(b)}} = b - 3a,\, b - 2a, \ovalbox{b - a}\, v :\, , b, \ovalbox{b - a}\ \langle...\rangle$ %
\href{\Scan{LH_37_05_196v}}{\lbrack196\:v\textsuperscript{o}\rbrack}\ %
\textit{(2)}~%
\bars{%
\textit{a} habet %
\lbrack...\rbrack\ velocitatem \textit{(v)}, dat %
\bars{ipsi \textit{b} \textit{erg. u. gestr.}} velocitatem \textit{z} ipsi \textit{b}. 
\lbrack...\rbrack\ velocitatem \textit{((v))}, \dots\dots\dots\ \textit{(b)} %
\textit{(1)~$\displaystyle\frac{\textit{(b)} - a}{\textit{(b)} + a} \textit{(v)} = \displaystyle\frac{b - 2a,\textit{(b)}-a}{b,\, a + b} v$} %
\textit{(2)}~$(v) = \displaystyle\frac{b - a}{b + a} v$ %
\lbrack...\rbrack\ seu $\textit{(b)}=b - 2a.$
\textit{erg.}%
} %
Ergo apparet series%
~\textit{L}%
}}
\lbrack%
\href{\Scan{LH_37_05_196v}}%
{196\,v\textsuperscript{o}}%
\rbrack%
\pend%
\vspace{0.5em}
\pstart%
\noindent%
\textit{a} habet velocitatem \textit{v},%
\indexS{velocitas corporis incurrentis}
\hspace{6mm}%
ictu retinet velocitatem \textit{(v)},%
\indexS{ictus}%
\indexS{velocitas retenta post ictum}
\hspace{6mm}%
dat velocitatem \textit{z} ipsi \textit{b}.%
\pend%
\pstart%
\noindent%
\textit{a} habet velocitatem \textit{(v)},
\hspace{3.3mm}%
ictu retinet velocitatem \textit{((v))},%
\indexS{ictus}%
\indexS{velocitas retenta post ictum}
\hspace{3.3mm}%
\dots\dots\dots\dots\dots\dots\dots
\textit{(b)}.
\pend%
\newpage%
\pstart%
$\textit{(v)} = \displaystyle\frac{b - a}{b + a} v$ %
\hspace{4mm}%
$z = \displaystyle\frac{2a}{a + b} v$
\pend
\vspace{1.25em}%
\pstart%
\noindent%
Fit
$z = \displaystyle\frac{2a}{a + \textit{(b)}}\textit{(v)} = \displaystyle\frac{2a}{a + b}v.$%
\hfill%
Sed
$\textit{(v)} = \displaystyle\frac{b - a}{b + a}v.$%
\hfill%
Ergo fit
$\displaystyle\frac{2a\cdot\overline{b - a}}{a + \textit{(b)},\, a + b} v = \displaystyle\frac{2a}{a + b}v.$
\newline%
Seu
$b - a = a + \textit{(b)}.$
\hspace{11mm}%
Seu
$\textit{(b)}=b - 2a.$%
\rule[0mm]{0mm}{4mm}%
\hspace{11mm}%
Ergo apparet series%
\indexS{series apparens}%
\lbrack:\rbrack%
\edlabel{LH_37_05_196-197_196r-v_uebergang-2}%
\pend%
\vspace{1.75em}%
\pstart%
$\textit{((v))}
=
\displaystyle\frac{\textit{(b)} - a}{\textit{(b)} + a}\textit{(v)}
=
\displaystyle\frac{b - 3a,\, b - a}{b - 2a + a,\, b + a}v
=
\displaystyle\frac{b - 3a}{b + a}v$%
\pend%
\vspace{0.5em}%
\pstart%
$\textit{((b))}
=
\textit{(b)} -2a = b - 4a$%
\pend%
\vspace{0.125em}%
\pstart%
$\textit{(((v)))}
=
\displaystyle\frac{\textit{((b))} - a}{\textit{((b))} + a} \textit{((v))}
=
\displaystyle\frac{b - 5a, \ovalbox{b - 3a} ,}{\ovalbox{b - 3a} , b + a} v$%
\pend%
\vspace{0.25em}%
\pstart%
$\textit{(((b)))}
=
\textit{((b))} - 2a = b - 6a$%
\pend%
\vspace{0.5em}%
\pstart%
$\textit{((((v))))}
=
\displaystyle\frac{\textit{(((b)))} - a}{\textit{(((b)))} + a} \textit{(((v)))}
=
\displaystyle\frac{b - 7a,\, b - 5a}{b - 5a,\, b + a} v$%
\pend%
\vspace{1.75em}%
\pstart%
\noindent%
Si \textit{z} sit data
fiet
$az + bz = 2av$
seu
$b : a = \overline{2v - z} : z.$%
\par%
\begin{edtabularr}%
\rule[0mm]{0pt}{9.5mm}%
\hspace*{6.5mm}%
\textit{a}%
&%
\phantom{[\textit{(((}}\textit{v}\phantom{\textit{)))}]}%
\edtext{}{%
\lemma{\textit{a}}%
\Bfootnote{%
\textit{(1)}~\textit{x} %
\textit{(2)}~\textit{v}%
~\textit{L}}}%
&%
$\textit{\phantom{(((}(v)\phantom{)))}} = \displaystyle\frac{b - a}{b + a}v\phantom{7}$%
&%
$z = \displaystyle\frac{2a}{a + b}v$%
\hspace{-6mm}%
&%
$\phantom{\textit{(((b)))} =} b \phantom{\ - 6a}$%
\hfill%
\\
\rule[0mm]{0pt}{5.5mm}%
$\bullet$%
&%
\phantom{[\textit{((}}\textit{(v)}\phantom{\textit{))}]}%
\edtext{}{%
\lemma{$\bullet$}%
\Bfootnote{%
\textit{(1)}~\textit{v} %
\textit{(2)}~\textit{(v)}%
~\textit{L}}}%
&%
$\textit{\phantom{((}((v))\phantom{))}} = \displaystyle\frac{b - 3a}{b + a}v$%
&%
&%
$\phantom{\textit{((}}\phantom{\textit{))}}\textit{(b)} = b - 2a$%
\hfill%
\\
\rule[0mm]{0pt}{5.5mm}%
$\bullet$%
&%
\phantom{[\textit{(}}\textit{((v))}\phantom{\textit{)}]}%
&%
$\textit{\phantom{(}(((v)))\phantom{)}} = \displaystyle\frac{b - 5a}{b + a}v$%
&%
&%
$\phantom{\textit{(}}\phantom{\textit{)}}\textit{((b))} = b - 4a$%
\hfill%
\\
\rule[-5,0mm]{0pt}{11.0mm}%
$\bullet$%
&%
[\textit{(((v)))}]%
\edtext{}{\lemma{}%
\Bfootnote{%
\textit{v}~\textit{L~ändert Hrsg.}}}%
&%
$\textit{((((v))))} = \displaystyle\frac{b - 7a}{b + a}v$%
&%
&%
$\textit{(((b)))} = b - 6a$%
\edlabel{LH_37_05_196-197_196v_irgh-1}%
\hfill%
\edtext{}{%
\xxref{LH_37_05_196-197_196v_irgh-1}{LH_37_05_196-197_196v_irgh-2}%
\lemma{$b - 6a$}%
\Bfootnote{%
\lbrack/\rbrack\ %
\textit{(1)}~Habemus ergo %
\textit{(2)}~Et sic %
\textit{(a)}~in infinitum
\textit{(b)}~porro, donec %
\lbrack...\rbrack\ ad \textit{b} ut
\textit{(aa)}~numerus %
\textit{(bb)}~unitas ad %
\lbrack...\rbrack\ omnis tandem %
\textit{(aaa)}~velocitas %
\textit{(bbb)}~vis corporis %
\lbrack...\rbrack\ massam aequivelocem %
\textit{(aaaa)}~. \lbrack/\rbrack\ \textlangle Ex his\textrangle\ %
\textit{(bbbb)}~praeditam velocitate \textit{z}. \lbrack/\rbrack\ Unde jam habemus%
~\textit{L}%
}}%
\end{edtabularr}%
\par%
\noindent%
Et sic porro,
\rule[0mm]{0pt}{4.5mm}%
donec exhauriatur \textit{b},
et proinde
si \textit{a} sit ad \textit{b} %
ut unitas%
\indexS{unitas}
ad numerum im\-pa\-rem%
\indexS{numerus impar}%
\lbrack,\rbrack\
omnis tandem vis corporis \textit{a}
translata erit in massam aequivelocem%
\indexS{massa aequivelox}%
\indexS{vis translata in massam aequivelocem}
praeditam velocitate~\textit{z}.%
\indexS{massa praedita velocitate}
Unde jam habemus%
\edlabel{LH_37_05_196-197_196v_irgh-2}
modum elegantem%
\indexS{modus elegans} %
ex legibus concursuum experientia comprobatis%
\indexS{lex concursus comprobata experientia}%
\indexS{experientia}
efficiendi
ut corporis
\edtext{\textit{a}}{%
\lemma{}%
\Bfootnote{%
\textit{a}%
~\textit{erg.\:L}%
}}
dati,
velocitate \textit{v} praediti%
\lbrack,\rbrack\
vis omnis transferatur
\edtext{%
in massam prius quiescentem%
\indexS{vis translata in massam quiescentem}%
\lbrack,\rbrack\
nunc omnes partes ae\lbrack q\rbrack uiveloces habentem.%
\indexS{massa habens partes aequiveloces}%
\indexS{pars aequivelox}
Quae cum sine ulla repugnantia%
\indexS{repugnantia}
possint connecti in unum corpus,%
\indexS{partes connexae in unum corpus}
si scilicet per reflexiones fiat%
\indexS{reflexio post ictum}
ut omnes accipiant directiones parallelas et in easdem partes,%
\indexS{directio accepta post ictum}%
\indexS{directio parallela}
hinc datur modus
omnem vim corporis minoris transferendi in corpus majus.%
\indexS{modus transferendi vim in corpus majus}%
\indexS{vis translata in corpus majus}%
}{%
\lemma{in}%
\Bfootnote{%
\textit{(1)}~corpus %
\textit{(2)}~massam quandam %
\textit{(a)}~aequi %
\textit{(b)}~aequivelocem. Sumatur corpus %
\textit{(3)}~it %
\textit{(4)}~alia %
\textit{(5)}~massam prius %
\lbrack...\rbrack\ omnes partes %
\textit{(a)}~omnes %
\textit{(b)}~%
\bars{aeviveloces \textit{ändert Hrsg.}} habentem %
\textit{(aa)}~,~seu ut omnis %
\textit{(aaa)}~motus %
\textit{(bbb)}~vis corporis minoris transferatur in majus %
\textit{(bb)}~.~Quae cum %
\lbrack...\rbrack\ per reflexiones %
\textit{(aaa)}~omnes habeant %
\textit{(bbb)}~fiat ut %
\textit{(aaaa)}~accipiant %
\textit{(bbbb)}~omnes accipiant %
\lbrack...\rbrack\ hinc datur %
\textit{(aaaaa)}~modum %
\textit{(bbbbb)}~modus %
\textit{(aaaaa-a)}~corporis minoris %
\textit{(bbbbb-b)}~omnem vim %
\lbrack...\rbrack\ corpus majus.%
~\textit{L}%
}}
Quanquam in rem nostram%
\indexS{res nostra}
parvum referat,
\edtext{\lbrack utrum\rbrack}{%
\lemma{}%
\Bfootnote{%
utrum \textit{erg. Hrsg.}}}
sint connexa an non
quae nunc aequivelociter moventur.
Sufficit enim
\edtext{nos habere}{%
\lemma{nos}%
\Bfootnote{%
\textit{(1)}~aestimare %
\textit{(2)}~habere%
~\textit{L}%
}}
materiam%
\indexS{materia affecta ubique eodem modo}
ubique eodem modo affectam.
\pend%
\pstart%
Modus hic%
\indexS{modus transferendi vim in corpus majus}%
\indexS{vis translata in corpus majus}
\edtext{%
est:
dato corpore \textit{a}%
}{%
\lemma{est:}%
\Bfootnote{%
\textit{(1)}~pro corpore \textit{a} %
\textit{(2)}~dato corpore \textit{a}%
~\textit{L}%
}}
motum \textit{v} habente
assumantur alia corpora quiescentia,%
\indexS{corpus quiescens ante ictum}
\textit{b} quod sit ad \textit{a} %
ut numerus impar,%
\indexS{numerus impar}
verbi
\edtext{%
gratia 5,
ad unitatem,%
\indexS{unitas}
seu \textit{b}
quod sit 5\textit{a},
et \textit{(b)}
quod sit 3\textit{a},
et \textit{((b))}
quod sit 1\textit{a}.
His positis%
}{%
\lemma{gratia}%
\Bfootnote{%
\textit{(1)}~7, %
\textit{(2)}~5, ad unitatem, %
\textit{(a)}~et \textit{(b)} %
\textit{(b)}~seu %
\textit{(aa)}~7\textit{a} %
\textit{(bb)}~quod sit 7\textit{a} %
\textit{(c)}~seu \textit{b} quod sit %
\textit{(aa)}~7\textit{a} %
\textit{(bb)}~5\textit{a}, et \textit{(b)} quod sit %
\textit{(aaa)}~5\textit{a},
\textit{(bbb)}~3\textit{a}, et \textit{((b))} quod sit %
\textit{(aaaa)}~3\textit{a} %
\textit{(bbbb)}~1\textit{a} %
\textit{(aaaaa)}~,~et \textit{(((b)))} quod sit 1\textit{a}, qui %
\textit{(bbbbb)}~.~His positis%
~\textit{L}%
}}
primum \textit{a}%
\indexS{corpus incurrens in aliud quiescens}
\edtext{%
velocitate \textit{v}%
\indexS{velocitas corporis incurrentis}%
}{%
\lemma{}%
\Bfootnote{%
velocitate \textit{v} \textit{erg.\:L}%
}}
incurrat
\edtext{%
in \textit{b}
quiescens%
\indexS{corpus quiescens ante ictum}
quintuplum ipsius,%
}{%
\lemma{in}%
\Bfootnote{%
\textit{(1)}~\textlangle\textendash\textrangle\ %
\textit{(2)}~\textit{b} quiescens %
\textit{(a)}~\textlangle triplum\textrangle\ %
\textit{(b)}~quintuplum ipsius,%
~\textit{L}}}
et dabit ei velocitatem \textit{z},
seu
$\displaystyle\frac{2a}{a + b}v$
\edtext{%
sive%
\rule[0.0ex]{0pt}{3.2ex}
$\displaystyle\frac{1}{3}v,$
reflecteturque velocitate%
\indexS{corpus reflexum post ictum}%
\indexS{velocitas corporis reflexi}
$\displaystyle\frac{b - a}{b + a}v$
seu $\displaystyle\frac{2}{3}v,$%
}{%
\lemma{sive $\displaystyle\frac{1}{3}v,$}%
\Bfootnote{%
\textit{(1)}~sibique retinebit velocitatem $\displaystyle\frac{3}{4}v$ %
\textit{(2)}~reflecteturque velocitate $\displaystyle\frac{b - a}{b + a}v$ seu $\displaystyle\frac{2}{3}v,$%
~\textit{L}%
}}
et cum progressus%
\indexS{progressus totalis}
antea fuerit \textit{av},%
\indexS{progressus ante incursum}
nunc est
$\displaystyle\frac{5}{3}av - \displaystyle\frac{2}{3}av,$%
\indexS{progressus post incursum}%
\rule[-0.4ex]{0pt}{3.8ex}
id est \textit{av} ut ante.
Et vis viva
erat ante \textit{avv},%
\indexS{vis viva ante incursum}
nunc est
$\displaystyle\frac{4}{9}avv + \displaystyle\frac{5}{9}avv = avv.$%
\indexS{vis viva post incursum}
Res ergo succedit.
\pend%
\vspace{0.25em}%
\pstart%
Corpus jam \textit{a} reflexum%
\indexS{corpus reflexum post ictum}%
\indexS{corpus incidens in aliud quiescens}
incidat in
\edtext{%
corpus \textit{(b)}
\lbrack it\rbrack idem
quiescens triplum%
\indexS{corpus quiescens ante ictum}%
}{%
\lemma{in corpus \textit{(b)}}%
\Bfootnote{%
\bars{ititidem \textit{ändert Hrsg.}} %
quiescens %
\textit{(1)}~quintuplum %
\textit{(2)}~triplum%
}}
ipsius \textit{a}
et dabit ei velocitatem eandem \textit{z}
sive
\edlabel{LH_37_05_196-197_196v-197r_uebergang-1}%
$\displaystyle\frac{1}{3}v,$%
\rule[-0.0ex]{0pt}{3.6ex}%
\edtext{}{%
\xxref{LH_37_05_196-197_196v-197r_uebergang-1}{LH_37_05_196-197_196v-197r_uebergang-2}%
\lemma{$\displaystyle\frac{1}{3}v,$ et}%
\Bfootnote{%
\textit{(1)}~ipsum \textit{a} reflectetur velocitate $\displaystyle\frac{1}{3}v.$ Habemus ergo progressum totalem %
\textit{(2)}~quae proinde erit contraria %
\bars{directionis \textit{ändert Hrsg.}} %
primae, sed ipsum %
\bars{corpus \textit{gestr.}} \textit{a} %
\bars{reflectendo \textit{ändert Hrsg.}} redibit ad %
\lbrack...\rbrack\ praedita velocitate $\displaystyle\frac{1}{3}v,$ %
\textit{(a)}~in quae trans %
\textit{(b)}~nempe %
\textit{(c)}~\textit{b}, \textit{(b)} et \textit{a}, %
\textit{(d)}~nempe \textit{b} ut %
\lbrack...\rbrack\ vis ipsius~\textit{a} %
\bars{praedita \textit{ändert Hrsg.}} velocitate~\textit{v}. Quodsi %
\lbrack...\rbrack\ nihil referret, %
\textit{(aa)}~cum ipsum %
\textit{(bb)}~quia hoc %
\lbrack...\rbrack\ progressum totalem, %
\bars{nam \textit{streicht Hrsg.}} %
habebimus nunc $\displaystyle\frac{b}{3}v$%
~\textit{L}%
}}
et quae proinde erit contraria%
\indexS{velocitas contraria directioni primae}
directio\lbrack ni\rbrack\
primae,
sed ipsum \textit{a}
refle\lbrack xum\rbrack\
redibit ad directionem primam
celeritate%
\indexS{celeritas corporis reflexi}
$\displaystyle\frac{1}{3}v.$
\pend%
\newpage%
\pstart%
Et ita nacti sumus tria corpora omnia aequivelocia%
\indexS{corpora aequivelocia}
seu praedita velocitate
$\displaystyle\frac{1}{3}v,$ %
nempe \textit{b} ut~5,
et \textit{(b)} ut~3,
et \textit{a} ut~1,
in quae translata erit vis ipsius \textit{a}%
\indexS{vis translata in corpora aequivelocia}
praedit\lbrack i\rbrack\
velocitate~\textit{v}.
Quodsi vellemus \textit{a} ad quietem%
\indexS{quies}
redigi,%
\indexS{corpus redactum ad quietem}
tantum debemus ponere incursum ejus directum%
\indexS{incursus directus in corpus aequale}
in aliud ei aequale \textit{((b))}
quod sit 1\textit{a}.
Sed hoc nihil referret,
quia hoc ab ipsomet non differret.
Habebimus autem rursus successum.%
\indexS{successus}
Nam
\lbrack%
\href{\Scan{LH_37_05_197r}}%
{197\,r\textsuperscript{o}}%
\rbrack\ %
computando progressum totalem,%
\indexS{progressus totalis}
habebimus nunc
$\displaystyle\frac{b}{3}v$%
\edlabel{LH_37_05_196-197_196v-197r_uebergang-2}
seu
$\displaystyle\frac{5}{3}av,$
una cum
$\displaystyle\frac{a}{3}v,$
demto $\displaystyle\frac{\textit{(b)}}{3}v$ %
seu demto \textit{v}.
Atque adeo progressus erit%
\indexS{progressus totalis}
$\displaystyle\frac{5 + 1 - 3}{3}av,$ %
id est \textit{av},
ut ante.
Sed et vis viva%
\indexS{vis viva in incursu}
manebit,
nam erit
$\displaystyle\frac{b}{9}avv + \displaystyle\frac{\textit{(b)}}{9}avv + \displaystyle\frac{1}{9}avv,$
seu
$\displaystyle\frac{5 + 3 + 1}{9}avv,$
id est \textit{avv}.
\pend%
\vspace{0.5em}%
\pstart%
Hinc jam vicissim
\edtext{omnibus inversis}{%
\lemma{}%
\Bfootnote{%
omnibus inversis%
~\textit{erg.\:L}%
}}
habetur modus%
\indexS{modus transferendi motum in massam minorem}
ex receptis experientiaque%
\indexS{experientia}
stabilitis
concursuum Legibus%
\indexS{lex concursus recepta}%
\indexS{lex concursus stabilita experientia}
omnem motum massae majoris
transferendi in massam minorem%
\indexS{massa mota}%
\indexS{motus tanslatus in massam minorem}%
\lbrack;\rbrack\
\edtext{%
\edlabel{LH_37_05_196-197_197r_Papinus-1}%
quod Dominus Papinus%
\indexN{\textls{Papin} (Papinus), Denis 1647\textendash1713}
fieri posse negabat.%
\edlabel{LH_37_05_196-197_197r_Papinus-2}%
}{%
\lemma{}%
\Bfootnote{%
quod Dominus %
\lbrack...\rbrack\ posse negabat%
~\textit{erg.\:L}%
}%
\lemma{quod \lbrack...\rbrack\ negabat}%
\Cfootnote{%
\textsc{D.\;Papin}, "Mechanicorum de viribus motricibus sententia" (\textit{AE}, %
1691, S.\;9\textendash11).\cite{00537}%
}%
}
Sint tria
\edtext{%
corpora sphaerica,%
\indexS{corpus sphaericum}
\textit{L} ut 1,%
}{%
\lemma{corpora}%
\Bfootnote{%
\textit{(1)}~\textit{a} ut 1, %
\textit{(2)}~sphaerica, \textit{L} ut 1,%
~\textit{L}%
}%
}
et \textit{M} ut 3,
et \textit{N} ut 5,
quae ferantur
\edtext{%
centris suis%
}{%
\lemma{}%
\Bfootnote{%
centris suis%
~\textit{erg.\:L}%
}}
in eadem linea
\edtext{%
\textit{AB}%
}{%
\lemma{}%
\Bfootnote{%
\textit{AB}%
~\textit{erg.\:L}%
}}
velocitate prorsus eadem,
quae sit ut 1,
sitque corpus \textit{L} medium
inter corpora \textit{M} et \textit{N},
ita quidem ut \textit{M} veniat ab \textit{A}
et tendat versus \textit{B}, %
sed \textit{L} et \textit{N} contra
\edtext{veniant a \textit{B}%
}{%
\lemma{veniant}%
\Bfootnote{%
\textit{(1)}~ab \textit{A} %
\textit{(2)}~a \textit{B}%
~\textit{L}%
}}
et tendant versus \textit{A},%
\edlabel{LH_37_05_196-197_197r_dugt-1}%
\edtext{}{%
\xxref{LH_37_05_196-197_197r_dugt-1}{LH_37_05_196-197_197r_dugt-2}%
\lemma{versus \textit{A},}%
\Bfootnote{%
\textit{(1)}~licet intervallo aliquo ex %
\textit{(2)}~ita tamen ut \textit{L} %
\textit{(a)}~procedat %
\textit{(b)}~praecedat, %
\textit{(aa)}~sed %
\textit{(bb)}~\textit{N} \bars{vero \textit{erg.}} ex intervallo sequatur. %
\textit{(aaa)}~His positis dico futu %
\textit{(bbb)}~Et \textit{L} prius incurrat in %
\textit{(ccc)}~His positis %
\lbrack...\rbrack\ corpora extrema %
\bars{\bars{\textit{L} \textit{ändert Hrsg.}} et \textit{N} \textit{erg.}} redigantur ad %
\textit{(aaaa)}~motu %
\textit{(bbbb)}~quietem %
\textit{(aaaaa)}~omnisque demum motus %
\textit{(bbbbb)}~et motus supersit in solo corpore %
\bars{\textit{M} \textit{ändert Hrsg.}}. \lbrack/\rbrack\ Nam quia %
\lbrack...\rbrack\ ante \textit{N} et %
\bars{in eadem linea \textit{erg.}} in occursum %
\lbrack...\rbrack\ ut corpus %
\textit{(aaaaa-a)}~\textit{L} %
\textit{(bbbbb-b)}~\textit{M} %
\bars{ut~3 \textit{gestr.}} redigatur ad quietem,%
~\textit{L}%
}}
ita tamen ut \textit{L} praecedat,%
\indexS{corpus praecedens}
\textit{N} vero ex intervallo%
\indexS{intervallum}
sequatur.%
\indexS{corpus sequens}
His positis dico
futurum esse ut corpora extrema
\lbrack\textit{M}\rbrack\ et \textit{N}
redigantur ad quietem%
\indexS{corpus redactum ad quietem}%
\indexS{quies}
et motus supersit%
\indexS{motus superstes}
in solo corpore \lbrack\textit{L}\rbrack.
\pend%
\begin{edfigure}[1.250em][1.00em]%
\includegraphics[width=0.84\textwidth]{%
gesamttex/edit_VIII-4/images/LH_37_05_196-197_d_197r.pdf}%
\edcaption[-0.50em]%
{%
\lbrack\textit{Fig.\:1}\rbrack%
\edtext{}{\lemma{\lbrack\textit{Fig.\:1}\rbrack}%
\Cfootnote[nonum]{%
\textit{L} (Bl.\;\href{\Scan{LH_37_05_197r}}{197\,r\textsuperscript{o}}) überliefert drei gestrichene, hier nicht wiedergegebene Entwürfe zum Diagramm.%
}}%
}%
\end{edfigure}%
\newpage%
\pstart%
Nam quia \textit{L} praecedit%
\indexS{corpus praecedens}
ante \textit{N} %
et in eadem linea in occursum%
\indexS{occursus corporum}
tendit ipsius \textit{M},
tandem ei occurret,%
\indexS{corpus occurrens alteri}
quo facto %
futurum erit
ut corpus \textit{M} redigatur ad quietem%
\indexS{quies},%
\indexS{corpus redactum ad quietem}%
\edlabel{LH_37_05_196-197_197r_dugt-2}
et \textit{L} reflectatur%
\indexS{corpus reflexum post occursum}
celeritate dupla ejus%
\indexS{celeritas corporis reflexi}
quam prius habebat seu ut 2.
Ita enim idem erit progressus
qui ante,%
\indexS{progressus totalis}
eademque
\edtext{vis.%
\indexS{vis totalis}
Nam}{%
\lemma{vis}%
\Bfootnote{%
\textit{(1)}~vis %
\textit{(2)}~viva %
\textit{(3)}~. Nam%
~\textit{L}%
}}
progressus%
\indexS{progressus totalis}
\edtext{%
antea a \textit{B}
versus \textit{A}
in duobus corporibus
\textit{L} et \textit{M} %
erat%
}{%
\lemma{antea}%
\Bfootnote{%
\textit{(1)}~erat ver %
\textit{(2)}~a \textit{B} versus %
\lbrack...\rbrack\ et \textit{M} erat%
~\textit{L}%
}}
$1L - 1M,$
seu
$1L - 3L$
seu
$-\,2L,$%
\indexS{progressus ante occursum}
sed nunc etiam est
$-\,2L$%
\indexS{progressus post occursum}
quia \textit{L} reflectitur%
\indexS{corpus reflexum post occursum}
versus
\edtext{%
\lbrack\textit{B}\rbrack%
}{%
\lemma{}%
\Bfootnote{%
\textit{A} \textit{L~ändert Hrsg.}%
}}
celeritate 2.%
\indexS{celeritas corporis reflexi}
Vis
\edtext{%
autem antea erat
$1^2 \cdot L + 1^2 \cdot M$%
}{%
\lemma{autem}%
\Bfootnote{%
\textit{(1)}~viva %
\textit{(2)}~antea %
\textit{(a)}~erit %
\textit{(b)}~erat %
\textit{(aa)}~$1L + 1L$ %
\textit{(bb)}~$1L + 1M$ %
\textit{(cc)}~$1^2 \cdot L + 1^2 \cdot M\,$%
~\textit{L}%
}}
seu $1L + 1M$
seu $1L + 3L$
seu 4\textit{L}.
Nunc vero est $2^2L$
seu
\edlabel{37_05_196-197_3a}%
etiam 4\textit{L}.%
\edtext{}{%
\xxref{37_05_196-197_3a}{37_05_196-197_3b}%
\lemma{etiam 4\textit{L}.}%
\Bfootnote{%
\textit{(1)}~Res %
\textit{(2)}~Et hoc jam sufficit ad %
\textit{(3)}~Res ergo %
\lbrack...\rbrack\ ad intentum,%
~\textit{L}%
}}%
\pend%
\pstart%
Res ergo succedit,
quod jam sufficeret ad intentum,%
\indexS{intentum}%
\edlabel{37_05_196-197_3b}
nam ita duorum corporum
\textit{L} ut~1
et~\textit{M} ut~3
vis tota redacta est%
\indexS{vis duorum corporum redacta in unum}%
\indexS{vis tota}
in solum corpus \textit{L}.
Sed si jam porro
\edtext{%
corpus \textit{L} ut~1
tendens%
}{%
\lemma{}%
\Bfootnote{%
corpus \textit{L} %
\bars{ut 1 \textit{erg.}} %
tendens%
~\textit{L}%
}}
celeritate 2%
\indexS{celeritas corporis occurrentis}
ab \textit{A} versus \textit{B}
occurrat%
\indexS{corpus occurrens alteri}
\edtext{%
corpori \textit{N} ut 5
a \textit{B} 
versus \textit{A}%
}{%
\lemma{corpori \textit{N}}%
\Bfootnote{\hspace{-0,625mm}%
\textit{(1)}~ab \textit{A} %
\textit{(2)}~ut 5 a \textit{B} versus \textit{A}%
~\textit{L}%
}}
tendenti
celeritate~1,%
\indexS{celeritas corporis occurrentis}
similiter concursu%
\indexS{concursus corporum}
corpus \textit{N}
redigetur ad quietem,%
\indexS{corpus redactum ad quietem}%
\indexS{quies}
et corpus \textit{L}%
\indexS{corpus reflexum post occursum}
\edtext{%
reflectetur
ut tendat
celeritate 3%
\indexS{celeritas corporis reflexi} %
\lbrack a \textit{B}\rbrack\
versus
\lbrack\textit{A}\rbrack.
Nam%
}{%
\lemma{reflectetur}%
\Bfootnote{%
\textit{(1)}~ab \textit{A} ver %
\textit{(2)}~ut tendat celeritate 3, %
\bars{ab \textit{A} \textit{ändert Hrsg.}} %
versus
\bars{\textit{B} \textit{ändert Hrsg.}}
. Nam%
~\textit{L}%
}}
ita idem
qui ante progressus,%
\indexS{progressus totalis}
eademque vis.%
\indexS{vis totalis}
Nam progressus corporum \textit{L} et \textit{N}
a \textit{B} versus \textit{A},
\edtext{%
ante concursum%
\indexS{progressus ante concursum}%
}{%
\lemma{}%
\Bfootnote{%
ante concursum%
~\textit{erg.\:L}%
}}
erat
$1N - 2L$
seu
$5L - 2L = 3L,$
et
\edtext{%
nunc post concursum%
\indexS{concursus corporum}%
\indexS{progressus post concursum}
est rursus 3\textit{L}.%
}{%
\lemma{nunc}\Bfootnote{%
\textit{(1)}~\textlangle est\textrangle\ \textit{L} %
\textit{(2)}~post concursum est rursus 3\textit{L}.%
~\textit{L}%
}}
Et vis ante concursum%
\indexS{vis ante concursum}
$1^2N + 2^2L$
seu
$5L + 4L$
seu 9\textit{L}.
Sed post concursum est
$3^2L,$%
\indexS{vis post concursum}
id est itidem 9\textit{L}.
Atque eodem modo pergendo%
\lbrack,\rbrack\
plurium adhuc aliorum corporum vires%
\indexS{vis plurium corporum collecta in unum}
in unum corpus \textit{L} colligi possent.%
\edlabel{LH_37_05_196-197_197r_gewp-1}%
\edtext{}{%
\xxref{LH_37_05_196-197_197r_gewp-1}{LH_37_05_196-197_197r_gewp-2}%
\lemma{possent.}%
\Bfootnote{%
\textit{(1)}~Et %
\textit{(2)}~Atque vis omnis corporis majoris %
\textit{(3)}~Et in universum %
\textit{(a)}~dati corporis %
\textit{(b)}~da %
\textit{(c)}~potentia massae %
\textit{(aa)}~magnitudinis datae %
\textit{(bb)}~quae sit %
\lbrack...\rbrack\ datae potest %
\textit{(aaa)}~reduci %
\textit{(bbb)}~colligi in %
\textit{(aaaa)}~corpus %
\textit{(aaaaa)}~datum minus %
\textit{(bbbbb)}~massam datam minorem %
\textit{(bbbb)}~corpus datum minus %
\textit{(aaaaa)}~vel potentia dispergi in eo mo %
\textit{(bbbbb)}~potentia %
\textit{(ccccc)}~. \lbrack/\rbrack\ Ex his etiam%
~\textit{L}%
}}
Et in universum potentia massae%
\indexS{potentia massae}
quae sit magnitudinis%
\indexS{magnitudo massae data}%
\indexS{massa magnitudinis datae}
et velocitatis datae%
\indexS{velocitas massae data}%
\indexS{massa velocitatis datae}%
\lbrack,\rbrack\
potest colligi in corpus datum minus.%
\indexS{potentia collecta in corpus minus}%
\pend%
\pstart%
Ex his etiam%
\edlabel{LH_37_05_196-197_197r_gewp-2}
porro sequitur
duorum corporum vires%
\indexS{vis duorum corporum translata in alia}
\edtext{posse transferri}{%
\lemma{posse}%
\Bfootnote{%
\textit{(1)}~reddi %
\textit{(2)}~transferri%
~\textit{L}%
}}
in corpora%
\indexS{vis translata in corpora velocitatis datae}
quae sint ejusdem velocitatis datae,
ut proinde facile possint aestimari%
\indexS{vis aestimata}%
\indexS{aestimatio virium}
cum sint ut corpora.%
\indexS{vires ut corpora}
Dato enim corpore \textit{a} %
et velocitate
\edtext{ejus \textit{v},
ac corpore}{%
\lemma{ejus \textit{v},}%
\Bfootnote{%
\textit{(1)}~et corpore %
\textit{(2)}~ac corpore%
~\textit{L}%
}}
\textit{l}
et velocitate ejus \textit{y},
sumamus velocitatem \textit{z}%
\indexS{velocitas ut mensura communis}
\edtext{%
talem,
ut sit communis mensura%
\indexS{mensura communis}%
}{%
\lemma{talem,}%
\Bfootnote{%
\textit{(1)}~ut %
\textit{(a)}~2\textit{v} itemque 2\textit{y} %
\textit{(b)}~metiatur %
\textit{(2)}~ut sit communis mensura%
~\textit{L}%
}}
ipsarum \textit{v} et \textit{y},
quas pone commensurabiles;%
\indexS{velocitas commensurabilis}
ita $2v - z$ erit ad \textit{z}
ut numerus impar ad unitatem,
et $2y - z$ erit similiter ad \textit{z}
ut numerus impar%
\indexS{numerus impar}
ad unitatem;%
\indexS{unitas}
inde jam determinari poterunt
\edtext{cor\lbrack p\rbrack ora}{%
\lemma{}%
\Bfootnote{%
corbora%
~\textit{L~ändert Hrsg.}}}
quiescentia \textit{b}, \textit{(b)}, \textit{((b))}%
\indexS{corpus quiescens ante incursum}
et corpora \textit{m}, \textit{(m)}, \textit{((m))}
talia
\edtext{%
ut ipsius \textit{a}%
\indexS{corpus incurrens in alia plura}
incursu%
\indexS{incursus corporis in alia plura}%
}{%
\lemma{ut}%
\Bfootnote{%
\textit{(1)}~illorum incu %
\textit{(2)}~ipsius \textit{a} incursu%
~\textit{L}%
}}
in illa
et ipsius \textit{l} incursu%
\indexS{incursus corporis in alia plura}
in haec,%
\indexS{corpus incurrens in alia plura}
successive consumantur velocitates%
\indexS{velocitas consumta successive}
ipsorum \textit{a} et \textit{l};
et corpora omnia \textit{b}, \textit{(b)}, \textit{((b))} etc. %
item \textit{m}, \textit{(m)}, \textit{((m))} etc.
sint ejusdem
\lbrack%
\href{\Scan{LH_37_05_197v}}%
{197\,v\textsuperscript{o}}%
\rbrack\ %
omnia et utraque velocitatis.%
\indexS{velocitas post incursum}
Unde erit vis ipsius \textit{a}
ad vim ipsius \textit{l},
ut est $b + \textit{(b)} + \textit{((b))}$ etc. %
ad $m + \textit{(m)} + \textit{((m))}$ etc.
\pend%
\vspace{4.0em}%
\pstart%
\noindent%
\centering%
*     *     *
\pend%
\selectlanguage{ngerman}
\frenchspacing%
\vspace{4.0em}%
\begin{ledgroup}%
\footnotesize%
\pstart%
\noindent%
\textls{Ausführliche Datierungsgründe:}%
\pagelabel{60359_dat}
Die titellose und undatierte Aufzeichnung N.\;\ref{60359} ist auf einem Bogen überliefert, dessen Wasserzeichen %
im Leibniz-Nachlass nach heutigem Kenntnisstand nur für den Zeitraum von 1690 bis März 1698 belegt ist.
Daher ist anzunehmen, dass N.\;\ref{60359} innerhalb dieses Zeitraums abgefasst wurde.
Der ausdrückliche Verweis auf Denis Papin\indexN{\textls{Papin} (Papinus), Denis 1647\textendash1713} im Text (\refpassage{LH_37_05_196-197_197r_Papinus-1}{LH_37_05_196-197_197r_Papinus-2}) liefert aber einen genaueren Terminus post quem für die Datierung:
Papin habe \glqq geleugnet\grqq, dass im Stoß die gesamte Kraft übertragen werden könne.
Diese Stellungnahme findet sich tatsächlich erstmals im Aufsatz "Mechanicorum de viribus motricibus sententia", den Papin im Heft der \textit{Acta eruditorum} vom Januar 1691 (S.\;6\textendash13, bes.\ 9\textendash11)\cite{00537} veröffentlicht hatte.
Die Aufzeichnung N.\;\ref{60359} muss folglich ab Januar 1691 verfasst worden sein.
Zudem könnte das Tempus in Leibnizens Erwähnung (\textit{negabat}) auf einen mehrjährigen zeitlichen Abstand zwischen Papins\indexN{\textls{Papin} (Papinus), Denis 1647\textendash1713} Aufsatz und der Abfassung von N.\;\ref{60359} hindeuten.
Ein genauer Terminus ante quem für die Datierung kann zum Redaktionsschluss jedoch nicht genannt werden.
Es ist aber unwahrscheinlich, dass die Aufzeichnung N.\;\ref{60359} nach dem Belegzeitraum des Wasserzeichens entstand, d.\,h. nach März 1698.
\pend%
\end{ledgroup}%
\selectlanguage{latin}%
\frenchspacing%
\normalsize%
